% !TeX root = ../papers/nc-p3.tex
\section{同型変換で得られる内部の三曲線の全点正則性}
\label{sec:rank-fifteen-regular-orbits}
\begin{theorem}\label{thm:rank-fifteen-regular-orbits}
$U=zk^2$, $V=zm^2$ とする。正の固定曲線の通常な階数15窓で
次の三因子のいずれかが零なら、両表示の全根は標準Koszul左右AS正則である。
各始点の経路の直和は四つの四次中心元の多項式環上で階数256の自由加群になる。
\[
\begin{aligned}
 F_1&=(U^2+14U+16V-7)^2-16U(3U-1)^2,\\
 F_2&=(U^2-2U+16V+9)^2-16U(U-3)^2,\\
 F_3&=(V-4(U+1))^2-64U .
\end{aligned}
\]
三式は $(U,V)$ 平面上で相異なる既約曲線を定める。
従来の座標端を含む倍率 $U-V=2$, $U+V=0$ および八次曲線の倍率 $V=(U-1)^2$ と合わせ、
前節の群の作用でこれらの条件に移る全ての通常窓が認証される。
\end{theorem}
\begin{proof}
前節の $a=1$, $b=(K-1)/2$, $c=a+b$ と五つの二次式を使う。
第一表示で $n^2=1$ は $a,e$ の一方が零になる元の枠であり、
$n^2=-1$ は $P(n)=0$ より $m^2=-k^2$ となり $b,f$ の一方が零になる元の枠である。
ここでいう枠の四文字は前節の係数配列の八量と区別する。
従って座標端の条件は係数配列で $16\gamma_4^4=\gamma_8^4$ と書ける。
192変換における $(F8,F4)$ は、第一ラベルが $8,1,9$ の順に、第二ラベルが
それぞれ $(4,5,12,13),(2,3,10,11),(6,7,14,15)$ の十二対に尽きる。
符号に用いる四乗根は四乗で消える。初めの四対では $n^2,N^2$ のいずれかが $\pm1$ であり、
和 $U-1$、積 $V$ を用いると元の二つの座標端水準になる。

次の四対では $4A^2,4B^2$ のいずれかが $\pm b^2$ である。この二根の和は $4c$、
積は $V$ なので条件は $b^4\pm4cb^2+V=0$ となる。
最後の四対では $b,c$ を交換した $c^4\pm4bc^2+V=0$ が得られる。
$K\mapsto-K$ の共役な二式ずつを掛けると $F_1,F_2$ になる。
八次曲線は $\xi=N^2/n^2$, $\xi^2+\xi+1=0$、すなわち
$(L^2+M^2)^2=L^2M^2$ である。各変換の $4,5$ の符号指数は同じ偶奇なので、
他の二対の条件も $(A^2+B^2)^2=A^2B^2$ と $(C^2+D^2)^2=C^2D^2$ になる。
従って $V=16c^2$ または $16b^2$ であり、その積の式が $F_3$ である。
各 $F_j$ は $V$ に関する二次式で、その判別式は非零平方の $U$ 倍であるから既約である。
例えば $U=4$ では順に $V=23/16,-25/16,36$ が各曲線上にあり、
これらでは両四重被覆の全根が座標端の外にある。

この条件で選んだ変換先が通常なら、座標端の全点正則性または
定理\ref{thm:rank-fifteen-octic-all}を適用できる。変換の十分性には二積の生成次元を仮定しない。
変換先が非通常のときも、前節の係数配列は定義され、元のテンソルに独立共役とスカラー倍で対応する。
以下、このテンソルから巡回する四つの隣接像を取る。$K,H$ を変換先の $rk,rm$ に相当する
$(r+1/r+w+1/w)/4,(r+1/r-w-1/w)/4$ とする。
復元式で $s,t$ は単元である。十五係数が非零なので $H\ne0$, $K\ne\pm1$ である。
通常条件の失敗には $K=0$ または $K+H,K-H$ の一方が $0,1,-1$ となることが必要である。
実際 $r+1/r=2(K+H)$, $w+1/w=2(K-H)$ で、$w=\pm r^{\pm1}$ のときは $KH=0$ となる。
これを各曲線に代入すると、検査すべき値は次に限られる。
\[
\begin{array}{c|l|l}
 \text{変換先} & H^2 & K\text{ の候補}\\\hline
 n^2=1 & K^2-2 & 0,\ 3/2,\ -3/2\\
 n^2=-1 & -K^2 & 4K^4+1=0\\
 \xi^2+\xi+1=0 & (K^2-1)^2 & 0,\ 2,\ -2,\ K^2\pm K-1=0
\end{array}
\]

\path{computations/verify_rank_fifteen_regular_orbits.py} はこれら全ての還元を
$\mathbb F_{17^{16}}$ で列挙する。最後の行では
$n^2=(K^2-1)(1\pm\sqrt{-3})/2$ である。各 $K,H,n$ に対し
$r+1/r,w+1/w$ から $r,w$ を求め、$s^2=r,t^2=w$ の全根を取る。
全て単数係数のモニックな方程式で整に選べ、逆元の分母 $2,s,t,n$ はこの還元で単元である。
符号変換で同じ剰余軌道に入るものは一代表にまとめる。
係数 $\gamma$ から直接
\[
 W_{ijku}=\boldsymbol1_{i\oplus k=j\oplus u}
  \sum_{v=0}^3 \gamma_{(i\oplus k)+4v}(-1)^{v\cdot k+\beta v\cdot u}
\]
を作り、各始点で隣接二文字像の基底を求める。全代表で各隣接像は六次元、
長さ4,5の経路は $35,56$ 次元、四次中心の文字次元は $(3,2,0,0)$ で、
四つの中心四次式による商は四始点とも長さ13まで定理\ref{thm:rank-fifteen-small-residues}の
次元列を与える。この計算は非通常倍率でのGram逆行列を使わない。

元の通常窓の経路と四次中心には一般点から来る次元の下界があり、
変換後の四文字テンソルも基底変更でその下界を保つ。上の剰余で隣接像の階数が
落ちないので四つの関係空間が持ち上がり、経路と中心の次元の一致から中心切断も持ち上がる。
よって前定理と同じ自由性・Koszul性の判定が適用される。
根交換と反対代数の対応によりこの認証は元の倍率の全根・両表示に及ぶ。
\end{proof}
